% META: {"id": "TNSI-BDD-CR-022", "chapitre": "TNSI-BASES-DE-DONNEES", "type_objet": "cours", "capacites_codes": ["C5", "C6", "C7"], "sources_inspiration": [], "mode_creation": "ex_nihilo", "status": "generated"}

\section{Interroger une base : SELECT, FROM, WHERE}

% ============================================================
% STRATE 1 — Composants d'une requete, SELECT / FROM
% ============================================================

\definition[1]{
Le langage \textbf{SQL} (\textit{Structured Query Language}) permet d'interroger et de manipuler une base relationnelle. Une requête d'interrogation combine des \textbf{clauses}, chacune introduite par un \textbf{mot clé} : \texttt{SELECT} (quels attributs afficher), \texttt{FROM} (dans quelle relation), et optionnellement \texttt{WHERE} (quelle condition sur les tuples).
}

\margeAppui{
\texttt{SELECT} correspond à une \textbf{projection} (choisir des colonnes) ; \texttt{WHERE} correspond à une \textbf{sélection} (choisir des lignes selon une condition). Ce sont deux opérations indépendantes : on peut projeter sans sélectionner, ou sélectionner sans projeter (avec \texttt{SELECT~*}).
}

\textbf{Exemple.} Base d'exemple filée sur tout le chapitre (bibliothèque) :
\begin{sql}
CREATE TABLE Livre(id INTEGER PRIMARY KEY, titre TEXT, auteur TEXT, annee INTEGER);
INSERT INTO Livre VALUES
  (1, 'Le Petit Prince', 'Saint-Exupery', 1943),
  (2, '1984', 'Orwell', 1949),
  (3, 'La Peste', 'Camus', 1947);
\end{sql}
Afficher uniquement les titres et auteurs de tous les livres :
\begin{sql}
SELECT titre, auteur FROM Livre;
\end{sql}

% BEGIN-VERIFY
% import sqlite3
% conn = sqlite3.connect(":memory:")
% conn.execute("CREATE TABLE Livre(id INTEGER PRIMARY KEY, titre TEXT, auteur TEXT, annee INTEGER)")
% conn.executemany("INSERT INTO Livre VALUES (?,?,?,?)", [
%     (1, "Le Petit Prince", "Saint-Exupery", 1943),
%     (2, "1984", "Orwell", 1949),
%     (3, "La Peste", "Camus", 1947),
% ])
% conn.commit()
% cur = conn.execute("SELECT titre, auteur FROM Livre")
% assert cur.fetchall() == [
%     ("Le Petit Prince", "Saint-Exupery"),
%     ("1984", "Orwell"),
%     ("La Peste", "Camus"),
% ]
% END-VERIFY

% ============================================================
% STRATE 2 — WHERE, operateurs et conditions composees
% ============================================================

\propriete[Filtrer avec WHERE]{
La clause \texttt{WHERE} ne conserve que les tuples vérifiant une \textbf{condition}. Les opérateurs de comparaison usuels (\texttt{=}, \texttt{<}, \texttt{>}, \texttt{<=}, \texttt{>=}, \texttt{<>} pour différent) se combinent avec \texttt{AND}, \texttt{OR}, \texttt{NOT} pour des conditions composées ; \texttt{LIKE} permet une recherche approchée sur du texte (motif avec le joker \texttt{\%}).
}

\textbf{Exemple.} Livres parus après 1945 :
\begin{sql}
SELECT titre FROM Livre WHERE annee > 1945;
\end{sql}
Livres d'Orwell ou de Camus :
\begin{sql}
SELECT titre FROM Livre WHERE auteur = 'Orwell' OR auteur = 'Camus';
\end{sql}

% BEGIN-VERIFY
% import sqlite3
% conn = sqlite3.connect(":memory:")
% conn.execute("CREATE TABLE Livre(id INTEGER PRIMARY KEY, titre TEXT, auteur TEXT, annee INTEGER)")
% conn.executemany("INSERT INTO Livre VALUES (?,?,?,?)", [
%     (1, "Le Petit Prince", "Saint-Exupery", 1943),
%     (2, "1984", "Orwell", 1949),
%     (3, "La Peste", "Camus", 1947),
% ])
% conn.commit()
% cur = conn.execute("SELECT titre FROM Livre WHERE annee > 1945")
% assert set(r[0] for r in cur.fetchall()) == {"1984", "La Peste"}
% cur = conn.execute("SELECT titre FROM Livre WHERE auteur = 'Orwell' OR auteur = 'Camus'")
% assert set(r[0] for r in cur.fetchall()) == {"1984", "La Peste"}
% END-VERIFY

\erreurFrequente{
\textbf{Oublier les guillemets autour d'une chaîne de caractères.} \texttt{WHERE auteur = Orwell} est une erreur : SQL interprète \texttt{Orwell} sans guillemets comme un nom de colonne inexistant. Il faut écrire \texttt{WHERE auteur = 'Orwell'}, avec des apostrophes simples autour du texte.
}
